\def\DeckTitle{Decompose the Task}\def\DeckDay{Monday Moments · AI Fluency}\def\DeckWeek{Week 4}\def\DeckSource{AI Fluency I · canonical Monday Moment package}\def\DeckQuestion{How do we make a large task inspectable before we rush into it?}\def\DeckPromise{Break a vague task into bounded stages with visible interfaces and checkpoints.}\def\DeckMove{Separate the whole goal from the stages, dependencies, and checks that make progress visible.}\def\DeckExample{Build a small course landing page: context → content → links → mobile check → publish decision.}\def\DeckEvidence{A staged plan, one interface/checkpoint per stage, and a reason to revise.}\def\DeckModel{A decomposition is a map from one whole task to smaller bounded stages. Each stage should have a purpose, an input, an output, and a checkpoint.}\def\DeckBoundary{More boxes do not make a plan safer. If stages hide dependencies or omit the user's goal, decomposition has created ceremony, not clarity.}\def\DeckApplication{Start with the student's goal, then name the content needed, the stable routes, the rendering check, and the decision point before publishing.}\def\DeckDo{Write the whole task first. Circle the first irreversible or uncertain step, then split only where a new boundary improves inspection.}\def\DeckSupports{Decomposition can expose interfaces, missing context, dependencies, risks, and verification points before implementation.}\def\DeckLimits{A plan is a hypothesis, not proof that the work will succeed; decomposition does not remove judgment, testing, or changing requirements.}\def\DeckWatch{A component inventory with no whole-task goal, or a “test it” step with no concrete case.}\def\DeckExit{Name one task you will decompose, its first checkpoint, and the observation that would make you revise the route.}\def\DeckStudy{Read the canonical `week_04_decompose_the_task` Monday Moment, student activity, instructor guide, rubric, and portfolio artifact before class.}\documentclass[aspectratio=169,11pt]{beamer}
\usepackage[T1]{fontenc}\usepackage{lmodern}\usepackage{hyperref}\usepackage{slides/week2_local_ai_workbench/commons-week2-theme}
\title{\DeckTitle}\subtitle{\DeckDay · Computing Commons · \DeckWeek}\author{\DeckSource}\date{}
\begin{document}\sloppy\setlength{\emergencystretch}{3em}\addtolength{\textheight}{8pt}\AtBeginEnvironment{frame}{\enlargethispage{4pt}}
\begin{frame}[plain]\begin{tikzpicture}[remember picture,overlay]\fill[ccnavy](current page.south west)rectangle(current page.north east);\fill[ccblue](current page.south west)rectangle([xshift=4.1cm]current page.south west);\draw[ccgold,line width=5pt]([xshift=1cm,yshift=1.5cm]current page.south west)--([xshift=14.5cm,yshift=1.5cm]current page.south west);\end{tikzpicture}\vspace{1.2cm}\color{white}\LARGE\textbf{\DeckTitle}\par\vspace{.35cm}\large\color{white!85}\DeckDay · \DeckWeek\par\vspace{.9cm}\normalsize\color{white!75}\DeckQuestion\par\vspace{.2cm}\color{white}\textbf{\DeckPromise}\end{frame}
\begin{frame}{Today's finish line}\begin{columns}[T]\column{.31\linewidth}\doit{See the core move.}\whatmeans{\DeckMove}\column{.31\linewidth}\doit{Try it on a bounded example.}\whatmeans{\DeckExample}\column{.31\linewidth}\doit{State evidence and limits.}\whatmeans{\DeckEvidence}\end{columns}\vfill\begin{center}\Large\textcolor{ccgold!80!black}{\textbf{What does this prove, and what does it not prove?}}\end{center}\vfill\ladder\end{frame}
\begin{frame}{The mental model}\begin{columns}[T]\column{.48\linewidth}\begin{beamercolorbox}[wd=\linewidth,sep=.6em]{block body}\textcolor{ccblue}{\textbf{The idea}}\par\DeckModel\end{beamercolorbox}\column{.48\linewidth}\begin{beamercolorbox}[wd=\linewidth,sep=.6em]{block body}\textcolor{ccgold!80!black}{\textbf{The boundary}}\par\DeckBoundary\end{beamercolorbox}\end{columns}\vfill\findyourself{Before accepting a conclusion, name the observation that would support it.}\end{frame}
\begin{frame}{A worked computing example}\doit{Apply the idea to this bounded case: \DeckExample}\vspace{.35cm}\begin{columns}[T]\column{.46\linewidth}\textbf{First pass}\begin{itemize}\item What is known now?\item Which step or assumption is visible?\item Where is the checkpoint?\end{itemize}\column{.46\linewidth}\begin{beamercolorbox}[sep=.7em,wd=\linewidth]{block body}\textcolor{ccgreen}{\textbf{Professional move}}\par\vspace{.2cm}\DeckApplication\end{beamercolorbox}\end{columns}\vfill\textcolor{ccblue}{\textbf{Label the application:}} This is an instructor/computing translation, not a claim about the primary source.\end{frame}
\begin{frame}{Ask before the answer}\begin{center}\Large\textcolor{ccnavy}{\textbf{\DeckQuestion}}\end{center}\vspace{.55cm}\begin{columns}[T]\column{.46\linewidth}\begin{beamercolorbox}[wd=\linewidth,sep=.6em]{block body}\textcolor{ccgold!80!black}{\textbf{Predict}}\par What would you expect if the first story were incomplete?\end{beamercolorbox}\column{.46\linewidth}\begin{beamercolorbox}[wd=\linewidth,sep=.6em]{block body}\textcolor{ccblue}{\textbf{Compare}}\par What alternative explanation deserves a fair test?\end{beamercolorbox}\end{columns}\vfill\textcolor{ccblue}{\textbf{PAIR TALK:}} Make one claim, one reason, and one uncertainty visible.\end{frame}
\begin{frame}{Do this: make the move visible}\doit{\DeckDo}\begin{enumerate}\item Write the task, decision, or claim in one sentence.\item Mark the relevant stages, evidence, comparison, or criteria.\item Choose one checkpoint before you act or explain.\item Say what would make you revise, stop, or continue.\end{enumerate}\vfill\verify{A good first attempt leaves a receipt another person can inspect.}\end{frame}
\begin{frame}{What the source supports — and its limits}\begin{columns}[T]\column{.47\linewidth}\textbf{Supports}\begin{itemize}\item \DeckSupports\item A clearer question and a more bounded next move.\end{itemize}\column{.47\linewidth}\begin{beamercolorbox}[sep=.75em,wd=\linewidth]{block body}\textcolor{ccgold!80!black}{\textbf{Does not prove}}\par\vspace{.2cm}\DeckLimits\end{beamercolorbox}\end{columns}\vfill\securityquestions\end{frame}
\begin{frame}{Recovery: when the first story fails}\begin{columns}[T]\column{.47\linewidth}\begin{beamercolorbox}[sep=.8em,wd=\linewidth]{block body}\centering\Huge\textbf{PAUSE}\par\vspace{.2cm}\normalsize A failed check is information about the route.\end{beamercolorbox}\column{.47\linewidth}\textbf{Use one bounded recovery}\begin{enumerate}\item Preserve the observation and context.\item Name the assumption or step to inspect.\item Make one smaller retry or ask for help.\item Do not turn uncertainty into a confident story.\end{enumerate}\end{columns}\vfill\textcolor{ccgreen}{\textbf{Watch for:}} \DeckWatch\end{frame}
\begin{frame}{Keep the receipt}\begin{columns}[T]\column{.47\linewidth}\doit{Record the smallest useful evidence.}\begin{itemize}\item Claim or decision\item Observation and context\item Reasoning / comparison / checkpoint\item Result, limit, and next action\end{itemize}\column{.47\linewidth}\begin{beamercolorbox}[wd=\linewidth,sep=.6em]{block body}\textcolor{ccgold!80!black}{\textbf{Exit move}}\par\DeckExit\end{beamercolorbox}\end{columns}\vfill\receipt{A concise receipt beats a confident story.}\end{frame}
\begin{frame}[plain]\begin{tikzpicture}[remember picture,overlay]\fill[ccnavy](current page.south west)rectangle(current page.north east);\fill[ccgreen](current page.south west)rectangle([xshift=5.1cm]current page.south west);\end{tikzpicture}\vspace{1.8cm}\color{white}\Huge\textbf{Your next move}\par\vspace{.55cm}\Large\color{white!85}\DeckExit\par\vspace{1cm}\color{ccgold}\Large\textbf{How will you know?}\end{frame}
\end{document}

